% Source-faithful assembly of Mumford, Abelian Varieties, Chapter I,
% Section 1 and the beginning of Section 2.  Running headers and pedagogical
% asides have been omitted; mathematical remarks and proofs are retained.

\section{Complex Tori}
\label{mumford-frag-I1}
\dcref{M:ch1:I.1}

We investigate a compact connected complex Lie group $X$ of dimension $g$.
Let $V$ be the tangent space to $X$ at the identity point $e$.  For every
$v\in V$ there is a unique holomorphic homomorphism
\[
 \phi_v:\mathbf C\longrightarrow X
\]
whose differential takes the unit tangent vector at $0$ to $v$, and
$\phi_v(t)=\exp(tv)$, where $\exp(v)=\phi_v(1)$.  The differential of
$\exp$ at $0$ is the identity.
\source{mumford-abelian-varieties:page-0012}

\begin{theorem}[Commutativity of a compact connected complex Lie group]
\label{mumford-frag-commutative}
\dcref{M:ch1:I.1}
Let $X$ be a compact connected complex Lie group.  Then $X$ is commutative.
\source{mumford-abelian-varieties:page-0012}
\end{theorem}
\begin{proof}
For $x\in X$, let $C_x(y)=xyx^{-1}$.  The differential $(dC_x)_e$ is an
automorphism of $V$, and $x\mapsto(dC_x)_e$ is a holomorphic map from $X$ to
the finite-dimensional vector space $\operatorname{End}(V)$.  It is constant,
so $(dC_x)_e=(dC_e)_e=1_V$.

For any homomorphism $T:X_1\to X_2$ of complex Lie groups,
\[
 T(\exp_{X_1}y)=\exp_{X_2}((dT)_e y).
\]
This follows from the uniqueness of the homomorphisms
$t\mapsto\exp_{X_i}(tv)$.  Thus
$C_x(\exp y)=\exp((dC_x)_e y)=\exp y$, so $\exp(V)$ is in the center.
Since $d\exp$ is the identity, the implicit function theorem shows that
$\exp$ maps a neighborhood of $0$ homeomorphically onto a neighborhood of
$e$.  As $X$ is connected, $\exp(V)$ generates $X$, and hence $X$ is
commutative.
\end{proof}

\begin{theorem}[Uniformization]
\label{mumford-frag-uniformization}
\dcref{M:ch1:I.1}
\uses{mumford-frag-commutative}
The exponential map $\exp:V\to X$ is a surjective homomorphism of complex
Lie groups with kernel a lattice $U$ in $V$, and induces an isomorphism
$V/U\simeq X$.
\source{mumford-abelian-varieties:page-0013}
\end{theorem}
\begin{proof}
For $x,y\in V$, commutativity shows that
$t\mapsto(\exp tx)(\exp ty)$ is a holomorphic homomorphism whose tangent at
$0$ is $x+y$.  Uniqueness gives
$(\exp tx)(\exp ty)=\exp t(x+y)$, so $\exp$ is a homomorphism.  Its image is
open and closed and contains $e$, hence is all of $X$.

The kernel $U$ is discrete because $\exp$ is injective on a neighborhood of
$0$.  The induced map $V/U\to X$ is holomorphic; its tangent map is an
isomorphism, so its inverse is holomorphic near the identity and everywhere
by translations.  Thus $V/U\simeq X$.  Since a discrete subgroup of a vector
space with compact quotient is a lattice, $U$ is a lattice.
\end{proof}

From now on write the group law additively and denote the exponential
homomorphism by $\pi:V\to X$.
\source{mumford-abelian-varieties:page-0014}

\begin{proposition}[Divisibility and torsion]
\label{mumford-frag-torsion}
\dcref{M:ch1:I.1}
\uses{mumford-frag-uniformization}
As an abstract group, $X$ is divisible.  If $n\neq0$ and $X_n$ is the subgroup
of elements annihilated by $n$, then
\[
 X_n\simeq(\mathbf Z/n\mathbf Z)^{2g}.
\]
\source{mumford-abelian-varieties:page-0014}
\end{proposition}
\begin{proof}
By uniformization, as a real Lie group $X\simeq(\mathbf R/\mathbf Z)^{2g}$,
which has these properties.
\end{proof}

\begin{theorem}[Integral cohomology]
\label{mumford-frag-integral-cohomology}
\dcref{M:ch1:I.1}
There are canonical isomorphisms
\[
 H^r(X,\mathbf Z)\simeq
 \{\text{alternating $r$-forms }U\times\cdots\times U\to\mathbf Z\}.
\]
\source{mumford-abelian-varieties:page-0014}
\end{theorem}
\begin{proof}
$(V,\pi)$ is the universal covering space, so $U=\pi^{-1}(0)=\pi_1(X,0)$.
The assertion is clear for $r=1$ from
$H^1(X,\mathbf Z)=\operatorname{Hom}(\pi_1(X),\mathbf Z)$.  It suffices to
show that cup product induces
\[
 \Lambda^rH^1(X,\mathbf Z)\xrightarrow{\sim}H^r(X,\mathbf Z).
 \tag{*}
\]
If $(*)$ holds for $X_1,X_2$ with finitely generated cohomologies, the
Kunneth formula gives it for $X_1\times X_2$:
\[
\begin{array}{ccc}
\Lambda^rH^1(X_1\times X_2,\mathbf Z)&\longrightarrow&H^r(X_1\times X_2,\mathbf Z)\\
\Vert&&\Vert\\
\Lambda^r[H^1(X_1,\mathbf Z)\oplus H^1(X_2,\mathbf Z)]&&\\
\Vert&&\Vert\\
\displaystyle\bigoplus_{p+q=r}[\Lambda^pH^1(X_1,\mathbf Z)\otimes\Lambda^qH^1(X_2,\mathbf Z)]&\xrightarrow{\sim}&
\displaystyle\bigoplus_{p+q=r}H^p(X_1,\mathbf Z)\otimes H^q(X_2,\mathbf Z).
\end{array}
\]
Here the torsion term vanishes because $(*)$ implies torsion-freeness.  Since
$X$ is a product of circles, $(*)$ holds.
\end{proof}

\begin{lemma}[Interior multiplication]
\label{mumford-frag-interior}
\dcref{M:ch1:I.1}
If $W$ is a complex vector space and $D\in\operatorname{Hom}_{\C}(W,\C)$,
then $D$ extends to $D\mathbin{\lrcorner}:\Lambda^pW\to\Lambda^{p-1}W$ by
\[
 D\mathbin{\lrcorner}(X_1\wedge\cdots\wedge X_p)=
 \sum_{k=1}^p(-1)^{p-k}D(X_k)X_1\wedge\cdots\wedge\widehat X_k\wedge\cdots\wedge X_p.
\]
If $DX_0=1$, then
$D\mathbin{\lrcorner}(\omega\wedge X_0)+(D\mathbin{\lrcorner}\omega)\wedge X_0=\omega$.
\source{mumford-abelian-varieties:page-0018}
\end{lemma}

\begin{theorem}[Dolbeault cohomology of a complex torus]
\label{mumford-frag-dolbeault}
\dcref{M:ch1:I.1}
Let $T=\operatorname{Hom}_{\mathbf C}(V,\mathbf C)$ and
$\overline T=\operatorname{Hom}_{\mathbf C\text{-antilinear}}(V,\mathbf C)$.
Then
\[
 H^q(X,\mathcal O_X)\simeq\Lambda^q\overline T,
 \qquad
 H^q(X,\Omega^p)\simeq\Lambda^pT\otimes\Lambda^q\overline T.
\]
\source{mumford-abelian-varieties:page-0015}
\end{theorem}
\begin{proof}
Translation extends each $\alpha\in\Lambda^pT$ to an invariant holomorphic
$p$-form $\omega_\alpha$, with
$(\omega_\alpha)_x=T_{-x}^*(\alpha)$.  Hence
\[
 \mathcal O_X\otimes_{\mathbf C}\Lambda^pT\xrightarrow{\sim}\Omega^p,
 \qquad
 H^q(X,\Omega^p)\simeq H^q(X,\mathcal O_X)\otimes\Lambda^pT.
\]
The Dolbeault resolution and fineness give
\[
 H^q(X,\mathcal O_X)\simeq
 \frac{\{\bar\partial\text{-closed $(0,q)$-forms}\}}
 {\bar\partial\{\text{$(0,q-1)$-forms}\}}.
\]
Let $\mathcal C=\mathcal C^{0,0}$ and
$\mathcal C^{p,q}$ be smooth forms of type $(p,q)$.  There is an isomorphism
\[
 \phi_{p,q}:\mathcal C\otimes_{\mathbf C}
 [\Lambda^pT\otimes\Lambda^q\overline T]\xrightarrow{\sim}\mathcal C^{p,q}.
\]
Invariant forms are closed: on pulling back to $V$, a degree-one invariant
form is $d\alpha$, so its exterior derivative is $d^2\alpha=0$, and the
general case follows by products.

Put $\Lambda=\bigoplus_q\Lambda^q\overline T$ and
$\mathfrak a=\Gamma(X,\mathcal C)$.  Then
\[
 \mathfrak a\otimes_{\mathbf C}\Lambda\xrightarrow{\sim}
 \Gamma(X,\mathcal C^{0,\bullet}),
 \qquad H^q(X,\mathcal O_X)\simeq H^q(\mathfrak a\otimes_{\mathbf C}\Lambda),
\]
where $\bar\partial(f\otimes\alpha)=\bar\partial f\wedge\alpha$.
The inclusion $\iota:\Lambda\to\mathfrak a\otimes\Lambda$ is a cohomology
isomorphism.  To prove this, normalize the Euclidean measure $\mu$ on $X$
to have volume $1$ and let $\mu:\mathfrak a\to\mathbf C$ be integration.
For any vector space $W$, $\mu\otimes1_W$ gives
\[
 \mu_A:\mathfrak a\otimes_{\mathbf C}A^\bullet\to A^\bullet,
 \qquad \mu_A\iota=\operatorname{Id}_{A^\bullet}.
\]
\source{mumford-abelian-varieties:page-0015}\source{mumford-abelian-varieties:page-0016}\source{mumford-abelian-varieties:page-0017}
\end{proof}

\begin{lemma}[Integration kills invariant derivatives]
\label{mumford-frag-integration-lemma}
\dcref{M:ch1:I.1}
For every $\omega\in\mathfrak a\otimes_{\mathbf C}A^\bullet$,
$\mu_A(\bar\partial\omega)=0$.
\source{mumford-abelian-varieties:page-0017}
\end{lemma}
\begin{proof}
By $A^\bullet$-linearity it is enough to consider $\bar\partial f$ for
$f\in\mathfrak a$.  In a basis $\omega_1,\ldots,\omega_n$ of $\overline T$,
write $\bar\partial f=\sum h_i\otimes\omega_i$.  Each $h_i$ is an invariant
directional derivative $D(f)$.  Pulling to $V$, periodicity and integration
over a fundamental domain give
\[
 \int_{V/U}D(f)\,dx=0,
\]
so every coefficient has zero integral.
\end{proof}

\begin{proof}[Proof (continued)]
For $\lambda\in U^*=\operatorname{Hom}(U,\mathbf Z)$, let
$e_\lambda\circ\pi(x)=e^{2\pi i\lambda(x)}$ and define Fourier coefficients
\[
 Q_\lambda(f)=\mu(e_{-\lambda}f),
 \qquad f=\sum_{\lambda\in U^*}e_\lambda\otimes Q_\lambda(f).
\]
Choosing a Hermitian norm and writing $\overline C(\lambda)$ for the
$\overline T$-projection of $\lambda$, one has
\[
 Q_\lambda(\bar\partial\omega)=(-1)^p2\pi i
 [Q_\lambda(\omega)\wedge\overline C(\lambda)].
\]
\source{mumford-abelian-varieties:page-0017}\source{mumford-abelian-varieties:page-0018}
\end{proof}

\begin{lemma}[Fourier coefficients and $\bar\partial$]
\label{mumford-frag-fourier-lemma}
\dcref{M:ch1:I.1}
For $\omega\in\mathfrak a\otimes_{\mathbf C}A^p$,
\[
 Q_\lambda(\bar\partial\omega)=(-1)^p2\pi i\,
 [Q_\lambda(\omega)\wedge\overline C(\lambda)].
\]
\source{mumford-abelian-varieties:page-0018}
\end{lemma}
\begin{proof}
Since $\bar\partial e_{-\lambda}=-2\pi i e_{-\lambda}\otimes
\overline C(\lambda)$, apply the preceding lemma to
$\bar\partial(\omega e_{-\lambda})$ and use the graded Leibniz rule.
The rapid-decay assertion is the standard Fourier-series isomorphism for
smooth functions on the compact torus.
\end{proof}

\paragraph{Fourier homotopy.}
For $\omega\in\mathfrak a\otimes\Lambda^p$, define $k(\omega)$ by
\[
 Q_\lambda(k\omega)=(-1)^{p-1}\lambda^*\mathbin{\lrcorner}Q_\lambda(\omega)
 \quad(\lambda\neq0),
 \qquad Q_0(k\omega)=0,
\]
where
$\lambda^*(x)=\langle x,\overline C(\lambda)\rangle/
(2\pi i\|\overline C(\lambda)\|^2)$.  Then
\[
 \bar\partial k+k\bar\partial=1_{\mathfrak a\otimes\Lambda^\bullet}-\iota\mu_A.
\tag{*}
\]
\source{mumford-abelian-varieties:page-0019}
\begin{proof}[Proof (continued)]
\textit{Proof.}
The Fourier coefficients have rapid decay, so exactly one such $k(\omega)$
exists.  For $\lambda\neq0$, Lemma~\ref{mumford-frag-fourier-lemma} and
interior multiplication give
\[
\begin{aligned}
 Q_\lambda(\bar\partial k\omega+k\bar\partial\omega)
 &=2\pi i\left[(-1)^{p-1}(\lambda^*\mathbin{\lrcorner}Q_\lambda\omega)
 \wedge\overline C(\lambda)+\lambda^*\mathbin{\lrcorner}
 (Q_\lambda\omega\wedge\overline C(\lambda))\right]\\
 &=Q_\lambda(\omega),
\end{aligned}
\]
and $Q_\lambda(\iota\mu_A\omega)=0$.  For $\lambda=0$, both homotopy terms
have zero coefficient and $Q_0(\omega)=Q_0(\iota\mu_A\omega)$.  Thus $(*)$
holds.  It implies that $\mu_A$ commutes with $\bar\partial$ and that
$\iota\mu_A$ is homotopic to the identity.
\end{proof}

\begin{remark}
\label{mumford-frag-cup-product}
\dcref{M:ch1:I.1}
Under $H^q(X,\mathcal O_X)\simeq\Lambda^q\overline T$, cup product is
exterior product.  This follows by computing cup product in the differential
graded algebra resolving $\mathcal O_X$.
\source{mumford-abelian-varieties:page-0019}
\end{remark}

\begin{corollary}
\label{mumford-frag-cup-cor}
\dcref{M:ch1:I.1}
The map induced by cup product
\[
 \Lambda^qH^1(X,\mathcal O_X)\longrightarrow H^q(X,\mathcal O_X)
\]
is an isomorphism.
\source{mumford-abelian-varieties:page-0019}\end{corollary}

\begin{remark}[de Rham and Hodge decomposition]
\label{mumford-frag-hodge}
\dcref{M:ch1:I.1}
The de Rham resolution gives
\[
 H^n(X,\C)\simeq
 \frac{\{d\text{-closed $n$-forms}\}}{d\{\text{$(n-1)$-forms}\}}.
\]
Every closed form differs by an exact form from a unique invariant form, so
\[
 H^n(X,\C)\simeq\Lambda^n\operatorname{Hom}_{\mathbf R}(V,\C).
\]
Since
\[
 \operatorname{Hom}_{\mathbf R}(V,\C)\simeq T\oplus\overline T,
\]
we obtain
\[
 H^n(X,\C)\simeq\Lambda^n(T\oplus\overline T)
 \simeq\bigoplus_{p+q=n}(\Lambda^pT\otimes\Lambda^q\overline T)
 \simeq\bigoplus_{p+q=n}H^q(X,\Omega^p).
\]
This is the Hodge decomposition.
\source{mumford-abelian-varieties:page-0020}\end{remark}

\begin{proposition}[Comparison of the three $H^1$ evaluations]
\label{mumford-frag-comparison}
\dcref{M:ch1:I.1}
Under the evaluations
\[
 H^1(X,\mathbf Z)\simeq U^*,\qquad
 H^1(X,\C)\simeq T\oplus\overline T,\qquad
 H^1(X,\mathcal O_X)\simeq\overline T,
\]
the maps induced by $\mathbf Z\to\C\to\mathcal O_X$ are the inclusion
$U^*\subset T\oplus\overline T$ and projection onto $\overline T$.
\source{mumford-abelian-varieties:page-0020}\source{mumford-abelian-varieties:page-0021}\source{mumford-abelian-varieties:page-0022}\source{mumford-abelian-varieties:page-0023}
\end{proposition}
\begin{proof}
For $\beta:H^1(X,\C)\to H^1(X,\mathcal O_X)$, compare the de Rham and
Dolbeault resolutions.  The projection
$C_{0,1}:\mathcal E^1=\mathcal E^{1,0}\oplus\mathcal E^{0,1}\to\mathcal E^{0,1}$
takes closed forms to $\bar\partial$-closed forms, giving a commutative
diagram of the two resolutions.  On cohomology this gives
\[
\begin{array}{ccccc}
T\oplus\overline T&\longrightarrow&H^0(X,\mathcal E^1_{d\text{-closed}})&\xrightarrow{\delta}&H^1(X,\C)\\
\downarrow\scriptstyle{\rm projection}&&\downarrow\scriptstyle{C_{0,1}}&&\downarrow\scriptstyle{\beta}\\
\overline T&\longrightarrow&H^0(X,\mathcal E^{0,1}_{\bar\partial\text{-closed}})&\xrightarrow{\delta}&H^1(X,\mathcal O_X).
\end{array}
\]
Thus $\beta$ is the expected map and is surjective.

For $\alpha:H^1(X,\mathbf Z)\to H^1(X,\C)$, let $a\in H^1(X,\mathbf Z)$.
For $u\in U$, use the loop $\phi_u(t)=\pi(tu)$ on $S^1=\mathbf R/\mathbf Z$
and the canonical isomorphism $\epsilon:H^1(S^1,\mathbf Z)\simeq\mathbf Z$:
\[
 \widetilde a(u)=\epsilon(\phi_u^*a).
\]
If $b\in T\oplus\overline T$ represents $\alpha(a)$, then
\[
 \widetilde a(u)=\epsilon(\phi_u^*a)
 =\epsilon(\delta(\phi_u^*\omega_b))
 =\int_{S^1}\phi_u^*(\omega_b)
 =\int_0^u\pi^*(\omega_b)=b(u).
\]
Indeed, if $\eta$ is a $1$-form on $S^1$ and $\delta(\eta)$ is its image in
$H^1(S^1,\C)$, then the canonical isomorphism
$\epsilon:H^1(S^1,\C)\simeq\C$ satisfies
\[
 \epsilon(\delta(\eta))=\int_{S^1}\eta.
\]
Thus $\widetilde a$ is the restriction of $b$ to $U$.  Compatibility with
cup products gives the corresponding commutative diagram in degree $n$:
\[
\begin{array}{ccccc}
H^n(X,\mathbf Z)&\xrightarrow{\sim}&\Lambda^nU^*\\
\downarrow&&\downarrow\Lambda^n(U^*\subset T\oplus\overline T)\\
H^n(X,\C)&\xrightarrow{\sim}&\displaystyle\bigoplus_{p+q=n}\Lambda^pT\otimes\Lambda^q\overline T\\
\downarrow&&\downarrow\\
H^n(X,\mathcal O_X)&\xrightarrow{\sim}&\Lambda^n\overline T.
\end{array}
\]
The right vertical map in the last row is projection onto the $p=0,q=n$
factor.
\end{proof}

\section{Line bundles on a complex torus}
\label{mumford-frag-I2}
\dcref{M:ch1:I.2}

\begin{theorem}[Vanishing on affine space]
\label{mumford-frag-affine-vanishing}
\dcref{M:ch1:I.2}
For every integer $p>0$, $H^p(\mathbf C^N,\mathcal O)=0$.
\source{mumford-abelian-varieties:page-0024}\end{theorem}

\begin{corollary}[Triviality on affine space]
\label{mumford-frag-affine-line-bundles}
\dcref{M:ch1:I.2}
For every $p>0$, $H^p(\mathbf C^N,\mathcal O^*)=0$.  In particular, every
holomorphic line bundle on $\mathbf C^N$ is trivial.
\source{mumford-abelian-varieties:page-0024}
\end{corollary}
\begin{proof}
Use
\[
 0\longrightarrow\mathbf Z\longrightarrow\mathcal O
 \xrightarrow{e^{2\pi i(\ )}}\mathcal O^*\longrightarrow0
\]
and $H^p(\mathbf C^N,\mathbf Z)=0$ for $p>0$.
\end{proof}

Let $L$ be a holomorphic line bundle on $X$.  Its pullback to $V$ is trivial;
choose $\chi:\pi^*L\xrightarrow{\sim}\mathbf C\times V$.  The canonical
$U$-action then has the form
\[
 (\alpha,z)\longmapsto\phi_u(\alpha,z)=(e_u(z)\alpha,z+u),
 \qquad e_u\in H^0(V,\mathcal O_V^*).
 \tag{A}
\]
The group law is equivalent to
\[
 e_{u+u'}(z)=e_u(z+u')e_{u'}(z).
\]
Changing $\chi$ by multiplication by $f\in H^0(V,\mathcal O_V^*)$ replaces
the cocycle by
\[
 e'_u(z)=e_u(z)f(z+u)f(z)^{-1}.
\]
\source{mumford-abelian-varieties:page-0024}\source{mumford-abelian-varieties:page-0025}

\begin{theorem}[Factors of automorphy]
\label{mumford-frag-factors}
\dcref{M:ch1:I.2}
Writing $H^*=H^0(V,\mathcal O_V^*)$, the above construction gives an
isomorphism
\[
 H^1(U,H^*)\xrightarrow{\sim}H^1(X,\mathcal O_X^*).
\]
Explicitly, a $1$-cocycle $\{e_u\}$ defines the quotient of
$\mathbf C\times V$ by (A), and every line bundle arises in this way.
More generally there is a natural map
\[
 H^i(U,\Gamma(\mathcal F,\pi^*\mathcal F))\longrightarrow H^i(X,\mathcal F),
\]
When $H^i(V,\pi^*\mathcal F)=0$ for $i\ge1$, this natural map is an
isomorphism; in the present case this applies because
$H^i(V,\mathcal O_V^*)=0$ for $i\ge1$.
\source{mumford-abelian-varieties:page-0025}
\end{theorem}
\begin{proof}
The cocycle construction in the forward direction was described above.  In
the reverse direction, quotienting by the displayed action gives a line
bundle, and cohomologous cocycles give isomorphic bundles.

To identify this with the Cech construction, choose a covering $\{V_i\}$ by
small connected open sets and connected lifts $W_i$.  Then
\[
 \pi^{-1}(V_i)=\coprod_{u\in U}(u+W_i),\qquad
 \pi_i:W_i\xrightarrow{\sim}V_i,
\]
and, when $V_i\cap V_j\ne\varnothing$, there is $u_{ij}\in U$ with
\[
 \pi_i^{-1}(V_i\cap V_j)=\pi_j^{-1}(V_i\cap V_j)+u_{ij}.
\]
The Cech cocycle is
\[
 f_{ij}(z)=e_{u_{ij}}(\pi_i^{-1}(z)).
\]
The associated line bundle is obtained by patching
$(\alpha,x)\mapsto(\alpha f_{ij}(x),x)$ over $V_i\cap V_j$.  Pulling this
description to $W_i$ gives
\[
 (\alpha,x)\longmapsto
 (\alpha f_{ij}(\pi_i(x)),x+u_{ij})=\phi_{u_{ij}}(\alpha,x),
\]
so it is precisely the quotient by $U$.
\end{proof}


On any complex analytic space the exponential sequence is exact:
\[
 0\longrightarrow\mathbf Z\longrightarrow\mathcal O_X
 \xrightarrow{e^{2\pi i(\ )}}\mathcal O_X^*\longrightarrow0.
\]
\footnote{With the action of $U$ on $H^*$ given by $(uh)(z)=h(z-u)$,
$u\in U$, $z\in V$, if $e_u$ satisfies the condition above then
$f_{u_{ij}}=e_{-u_{ij}}$ is a $1$-cocycle for this action, and conversely.
Thus the associated $1$-cocycle is given by
\[
 f_{ij}(z)=f_{-u_{ij}}(\pi_i^{-1}(z))
 =e_{u_{ij}}(\pi_i^{-1}(z)).
\]}
\source{mumford-abelian-varieties:page-0026}
